\documentclass[a4paper]{article}

\usepackage{graphicx}
\usepackage{amsmath,amssymb}
\usepackage{minted}
\usepackage{multirow}
\usepackage{float}
\usepackage{todonotes}
\usepackage{forest}
\usepackage{xstring}
\usepackage{xcolor}
\usepackage[most]{tcolorbox}
\usepackage{xparse}
\usepackage{natbib}

\usetikzlibrary{graphs,graphdrawing}
\usegdlibrary{layered}

% Benötigt: \usepackage{xcolor}
\definecolor{TranslationGray}{RGB}{120,120,120}
\newcommand{\EN}[1]{\par{\small\textcolor{TranslationGray}{#1}}\vspace{0.4em}}

% for external graphviz
\input{/home/donkarlo/Dropbox/repo/nd_ai_project/src/nd_ai/knoweldge_representation/language/formal/computer/programming/tex/latex/code_clips/external_graphviz.tex}

% for tags
\input{/home/donkarlo/Dropbox/repo/nd_ai_project/src/nd_ai/knoweldge_representation/language/formal/computer/programming/tex/latex/parts/control_sequence/kind/semtag/semtag.tex}

% for links
\input{/home/donkarlo/Dropbox/repo/nd_ai_project/src/nd_ai/knoweldge_representation/language/formal/computer/programming/tex/latex/code_clips/seehere.tex}

\title{Kasus}
\date{}

\begin{document}
    \maketitle
    
    
    \section{Kasus bei deutschen Verben}
        
        \textbf{Im Deutschen reicht es nicht immer, nur direktes und indirektes Objekt zu erkennen.}
        \EN{In German, it is not always enough to recognize only direct and indirect objects.}
        
        \textbf{Oft bestimmt das Verb selbst, welchen Kasus das Objekt bekommt.}
        \EN{Often the verb itself determines which case the object takes.}
        
        \textbf{Deshalb sollte man viele Verben zusammen mit ihrem Kasus lernen.}
        \EN{Therefore, many verbs should be learned together with their case.}
        
        \subsection{Verben mit Dativ}
            
            \textbf{Einige Verben brauchen immer oder meistens den Dativ.}
            \EN{Some verbs always or usually require the dative.}
            
            \textbf{helfen + Dativ}
            \EN{to help + dative}
            
            \textbf{Ich helfe dir.}
            \EN{I help you.}
            
            \textbf{danken + Dativ}
            \EN{to thank + dative}
            
            \textbf{Ich danke dem Großvater.}
            \EN{I thank the grandfather.}
            
            \textbf{antworten + Dativ}
            \EN{to answer + dative}
            
            \textbf{Ich antworte der Lehrerin.}
            \EN{I answer the teacher.}
            
            \textbf{gehören + Dativ}
            \EN{to belong to + dative}
            
            \textbf{Das Buch gehört mir.}
            \EN{The book belongs to me.}
        
        \subsection{Verben mit Akkusativ}
            
            \textbf{Die meisten Verben mit Objekt brauchen den Akkusativ.}
            \EN{Most verbs with an object require the accusative.}
            
            \textbf{sehen + Akkusativ}
            \EN{to see + accusative}
            
            \textbf{Ich sehe dich.}
            \EN{I see you.}
            
            \textbf{haben + Akkusativ}
            \EN{to have + accusative}
            
            \textbf{Ich habe einen Kaffee.}
            \EN{I have a coffee.}
            
            \textbf{kaufen + Akkusativ}
            \EN{to buy + accusative}
            
            \textbf{Ich kaufe das Brot.}
            \EN{I buy the bread.}
        
        \subsection{Verben mit zwei Objekten}
            
            \textbf{Wenn ein Satz zwei Objekte hat, ist die Person oft Dativ und die Sache Akkusativ.}
            \EN{When a sentence has two objects, the person is often dative and the thing is accusative.}
            
            \textbf{Ich gebe dir das Buch.}
            \EN{I give you the book.}
            
            \textbf{dir = Dativ}
            \EN{to you = dative}
            
            \textbf{das Buch = Akkusativ}
            \EN{the book = accusative}
            
            \textbf{Ich werfe dir den Ball zu.}
            \EN{I throw the ball to you.}
            
            \textbf{dir = Dativ}
            \EN{to you / receiver = dative}
            
            \textbf{den Ball = Akkusativ}
            \EN{the ball = accusative}
        
        \subsection{Wichtige Regel}
            
            \textbf{Für Deutsch muss man zuerst lernen, welchen Kasus ein Verb verlangt.}
            \EN{For German, one must first learn which case a verb requires.}
            
            \textbf{Direktes und indirektes Objekt helfen beim Verstehen, aber sie reichen nicht immer aus.}
            \EN{Direct and indirect objects help with understanding, but they are not always enough.}
            
            \textbf{Die beste Methode ist: Verb + Kasus + Beispiel lernen.}
            \EN{The best method is: learn verb + case + example.}
            
            \textbf{helfen + Dativ: Ich helfe dir.}
            \EN{to help + dative: I help you.}
            
            \textbf{sehen + Akkusativ: Ich sehe dich.}
            \EN{to see + accusative: I see you.}
            
            \textbf{geben + Dativ + Akkusativ: Ich gebe dir das Buch.}
            \EN{to give + dative + accusative: I give you the book.}
    
    
    \section{Kasus bei Präpositionen}
        
        \textbf{Nicht nur Verben bestimmen den Kasus. Auch Präpositionen bestimmen oft den Kasus.}
        \EN{Not only verbs determine the case. Prepositions often determine the case too.}
        
        \textbf{Eine Präposition ist ein kleines Wort wie mit, für, seit, wegen oder in.}
        \EN{A preposition is a small word like with, for, since, because of, or in.}
        
        \textbf{Nach manchen Präpositionen kommt immer Dativ.}
        \EN{After some prepositions, the dative always comes.}
        
        \textbf{Nach manchen Präpositionen kommt immer Akkusativ.}
        \EN{After some prepositions, the accusative always comes.}
        
        \textbf{Manche Präpositionen können Dativ oder Akkusativ haben.}
        \EN{Some prepositions can take dative or accusative.}
        
        \subsection{Präpositionen mit Dativ}
            
            \textbf{Diese Präpositionen brauchen immer den Dativ: aus, bei, mit, nach, seit, von, zu.}
            \EN{These prepositions always require the dative: from/out of, at/near, with, after/to, since/for, from/of, to.}
            
            \textbf{mit + Dativ}
            \EN{with + dative}
            
            \textbf{Ich spreche mit dem Lehrer.}
            \EN{I speak with the teacher.}
            
            \textbf{bei + Dativ}
            \EN{at / near / with + dative}
            
            \textbf{Ich bin bei meiner Mutter.}
            \EN{I am at my mother's place.}
            
            \textbf{nach + Dativ}
            \EN{after / to + dative}
            
            \textbf{Nach dem Kurs gehe ich nach Hause.}
            \EN{After the course I go home.}
            
            \textbf{von + Dativ}
            \EN{from / of + dative}
            
            \textbf{Das ist ein Brief von meinem Bruder.}
            \EN{That is a letter from my brother.}
            
            \textbf{zu + Dativ}
            \EN{to + dative}
            
            \textbf{Ich gehe zu meinem Professor.}
            \EN{I go to my professor.}
        
        \subsection{Die Präposition seit}
            
            \textbf{Seit ist eine Präposition mit Dativ.}
            \EN{Seit is a preposition with the dative.}
            
            \textbf{Man benutzt seit für eine Zeit, die in der Vergangenheit begonnen hat und bis jetzt dauert.}
            \EN{Seit is used for a time that began in the past and continues until now.}
            
            \textbf{seit + Dativ}
            \EN{since / for + dative}
            
            \textbf{Ich lerne seit einem Jahr Deutsch.}
            \EN{I have been learning German for one year.}
            
            \textbf{Er wohnt seit drei Monaten in Wien.}
            \EN{He has been living in Vienna for three months.}
            
            \textbf{Wir kennen uns seit letzter Woche.}
            \EN{We have known each other since last week.}
            
            \textbf{Ich warte seit zehn Minuten.}
            \EN{I have been waiting for ten minutes.}
            
            \textbf{Ich lebe seit 2024 in Österreich.}
            \EN{I have been living in Austria since 2024.}
        
        \subsection{Wichtige Formen mit seit}
            
            \textbf{seit einem Tag}
            \EN{for one day}
            
            \textbf{seit einer Woche}
            \EN{for one week}
            
            \textbf{seit einem Monat}
            \EN{for one month}
            
            \textbf{seit einem Jahr}
            \EN{for one year}
            
            \textbf{seit zwei Tagen}
            \EN{for two days}
            
            \textbf{seit drei Wochen}
            \EN{for three weeks}
            
            \textbf{seit vier Monaten}
            \EN{for four months}
            
            \textbf{seit fünf Jahren}
            \EN{for five years}
            
            \textbf{seit dem Morgen}
            \EN{since the morning}
            
            \textbf{seit der Kindheit}
            \EN{since childhood}
        
        \subsection{Seit mit Präsens}
            
            \textbf{Im Deutschen benutzt man mit seit oft das Präsens, wenn die Situation noch jetzt gilt.}
            \EN{In German, with seit one often uses the present tense when the situation is still true now.}
            
            \textbf{Ich wohne seit einem Jahr in Paris.}
            \EN{I have been living in Paris for one year.}
            
            \textbf{Ich arbeite seit zwei Stunden.}
            \EN{I have been working for two hours.}
            
            \textbf{Sie wartet seit zehn Minuten.}
            \EN{She has been waiting for ten minutes.}
            
            \textbf{Wir lernen seit drei Monaten Deutsch.}
            \EN{We have been learning German for three months.}
        
        \subsection{Häufiger Fehler mit seit}
            
            \textbf{Falsch: Ich freue mich, dass du hast seit ein Jahr in Paris leben.}
            \EN{Wrong: I am happy that you have since one year live in Paris.}
            
            \textbf{Richtig: Ich freue mich, dass du seit einem Jahr in Paris lebst.}
            \EN{Correct: I am happy that you have been living in Paris for one year.}
            
            \textbf{Nach dass steht das konjugierte Verb am Ende.}
            \EN{After dass, the conjugated verb goes to the end.}
            
            \textbf{Deshalb sagt man: dass du ... lebst.}
            \EN{Therefore one says: that you ... live.}
            
            \textbf{Nach seit kommt Dativ.}
            \EN{After seit comes dative.}
            
            \textbf{Deshalb sagt man: seit einem Jahr.}
            \EN{Therefore one says: for one year.}
            
            \textbf{Mit seit benutzt man oft Präsens, wenn die Handlung noch dauert.}
            \EN{With seit one often uses the present tense when the action is still continuing.}
        
        \subsection{Seit und seid}
            
            \textbf{Seit und seid klingen gleich, aber sie haben verschiedene Bedeutungen.}
            \EN{Seit and seid sound the same, but they have different meanings.}
            
            \textbf{seit = Präposition für Zeit}
            \EN{seit = preposition for time}
            
            \textbf{Ich lerne seit einem Jahr Deutsch.}
            \EN{I have been learning German for one year.}
            
            \textbf{seid = Verbform von sein}
            \EN{seid = verb form of to be}
            
            \textbf{Ihr seid sehr freundlich.}
            \EN{You are very friendly.}
            
            \textbf{Ihr seid seit einer Stunde hier.}
            \EN{You have been here for one hour.}
    
    
    \section{Zusammenfassung}
        
        \textbf{Bei deutschen Verben muss man oft lernen, welchen Kasus das Verb verlangt.}
        \EN{With German verbs, one often has to learn which case the verb requires.}
        
        \textbf{Bei deutschen Präpositionen muss man ebenfalls lernen, welchen Kasus die Präposition verlangt.}
        \EN{With German prepositions, one also has to learn which case the preposition requires.}
        
        \textbf{helfen + Dativ: Ich helfe dir.}
        \EN{to help + dative: I help you.}
        
        \textbf{sehen + Akkusativ: Ich sehe dich.}
        \EN{to see + accusative: I see you.}
        
        \textbf{geben + Dativ + Akkusativ: Ich gebe dir das Buch.}
        \EN{to give + dative + accusative: I give you the book.}
        
        \textbf{seit + Dativ: Ich lerne seit einem Jahr Deutsch.}
        \EN{since / for + dative: I have been learning German for one year.}
        
        \textbf{Die beste Methode ist: Wort + Kasus + Beispiel lernen.}
        \EN{The best method is: learn word + case + example.}
        
%        \nocite{*}
%        \bibliographystyle{plainnat}
%        \bibliography{/home/donkarlo/Dropbox/repo/nd_shared_research/bibliography/bibliography.bib}

\end{document}